\documentclass[12pt,a4paper]{article}
\usepackage[utf8]{inputenc}
\usepackage{amsmath,amssymb}
\usepackage{graphicx}
\usepackage[left=1.5cm,right=1cm,top=1.3cm,bottom=1.3cm]{geometry}
\usepackage{lastpage}
% Font Times New Roman (cần XeLaTeX + fontspec — uncomment nếu VPS hỗ trợ):
% \usepackage{fontspec}
% \setmainfont{Times New Roman}

\begin{document}

\begin{center}
\begin{tabular}{p{0.4\textwidth}p{0.6\textwidth}}
\textbf{SỞ GDĐT SƠN LA} & ĐỀ KIỂM TRA GIỮA KÌ 1 NĂM 2026-2027 \tabularnewline
\textbf{TRƯỜNG THPT CÂY MƠ} & \textbf{Môn: TOÁN 12} \tabularnewline
\textbf{Đề chính thức} & \textit{Thời gian: 90 phút, không kể thời gian phát đề} \tabularnewline
\textit{(Đề thi gồm có .... trang)} & \tabularnewline
Họ và tên: ...................................................................................... \newline Số báo danh: .................................................................................. & \textbf{Mã đề: 8618} \tabularnewline
\end{tabular}
\end{center}
\vspace{0.5cm}

\begin{center}
{\Large\textbf{ĐÁP ÁN VÀ LỜI GIẢI}}
\end{center}

\noindent\textbf{PHẦN I. Câu trắc nghiệm nhiều phương án lựa chọn.}

\medskip

\noindent\textcolor{blue}{\textbf{Câu 1.}}

\noindent \textbf{Đáp án D.}

\noindent\textbf{Lời giải.}

Trên $[1;6]$ ta có $f'(x)=0\Leftrightarrow 3x^2-39=0\Leftrightarrow x=\sqrt{13}$ hoặc $x=-\sqrt{13}\not\in[1;6]$.\\
$f(x)$ liên tục và xác định trên $[1;6]$, đồng thời $\begin{cases} f(1)=-38 \\f\left (\sqrt{13}\right )=-26\sqrt{13}\\f(6)=-18. \end{cases}$\\

	Vậy $\min\limits_{[1;6]}f(x)=f\left (\sqrt{13}\right )=-26\sqrt{13}$.

\medskip
\noindent\textcolor{blue}{\textbf{Câu 2.}}

\noindent \textbf{Đáp án C.}

\noindent\textbf{Lời giải.}

Hàm số đã cho liên tục trên đoạn $[2;5]$.\\

		Ta có $f'(x)=\dfrac{2\cdot(-1)-1\cdot 3}{(x-1)^2}=\dfrac{-5}{(x-1)^2}<0$ với mọi $x\in [2;5]$.\\

		Hàm số nghịch biến trên đoạn $[2;5]$ nên giá trị lớn nhất của hàm số trên đoạn này đạt được tại $x=2$.\\

		Giá trị lớn nhất là $\max\limits_{[2;5]} f(x)=f(2)=\dfrac{2\cdot 2+3}{2-1}=7$.

\medskip
\noindent\textcolor{blue}{\textbf{Câu 3.}}

\noindent \textbf{Đáp án C.}

\noindent\textbf{Lời giải.}

Hàm số bậc hai $y=x^2+2x-3$ có hệ số $a>0$ nên có giá trị nhỏ nhất trên tập xác định của nó.

\medskip
\noindent\textbf{PHẦN II. Câu trắc nghiệm đúng sai.}

\medskip

\noindent\textcolor{blue}{\textbf{Câu 1.}}

\noindent \textbf{Đáp án: ĐSĐĐ}

\noindent\textbf{Lời giải.}

a)  Với $x=4$ thì $f(4)=\dfrac{4^2-4-12}{4-2}=0$.

			Tập xác định $\mathscr{D}=\mathrm{R}\setminus \{2\}$.\\

			Ta có $\lim\limits_{x\to+\infty} \left[f(x)-(x+1)\right]=\lim\limits_{x\to+\infty} \left[x+1-\dfrac{10}{x-2}-(x+1)\right]=\lim\limits_{x\to+\infty}-\dfrac{10}{x-2}=0$.\\
$\lim\limits_{x\to-\infty} \left[f(x)-(x+1)\right]=\lim\limits_{x\to-\infty} \left[x+1-\dfrac{10}{x-2}-(x+1)\right]=\lim\limits_{x\to-\infty}-\dfrac{10}{x-2}=0$.\\

			Do đó đồ thị hàm số $f(x)$ có đường tiệm cận xiên là đường thẳng $y=x+1$.

			Do đó $x=2$ là tiệm cận đứng của đồ thị hàm số $y=f(x)=\dfrac{x^2-x-12}{x-2}$.\\

			Giao điểm của tiệm cận đứng và tiệm cận xiên là $(2;3)$.\\
Do đó đồ thị hàm số $f(x)$ có tâm đối xứng là điểm $I(2; 3)$.

			Ta có $f'(x)=\dfrac{(x^2-4x+14)}{(x-2)^2}$.\\

			Do đó $f'(3)=11$.\\

			Vậy phương trình tiếp tuyến $y=11(x-3)-6=11x-39$.\\

			Suy ra $m=11$, $n=-39$.\\

			Vậy $m-n=50$.

\medskip
\noindent\textbf{PHẦN III. Câu trả lời ngắn.}

\medskip

\noindent\textcolor{blue}{\textbf{Câu 1.}}

\noindent \textbf{Đáp số: 1}

\noindent\textbf{Lời giải.}

Dựa vào đồ thị của hàm số $y=f'(x)$ ta thấy $f'(x)=0\Leftrightarrow \left[\begin{array}{l}x=-1 \\x=1 \\x=4.\end{array}\right.$\\

		Ta có bảng biến thiên của hàm số $y=f(x)$
% [Hình TikZ không compile được: tikz_dae45c23a41f]

		Trên đoạn $[-1;4]$, hàm số $y=f(x)$ đạt giá trị lớn nhất tại $x=1$.

\medskip
\noindent\textcolor{blue}{\textbf{Câu 2.}}

\noindent \textbf{Đáp số: 4,5}

\noindent\textbf{Lời giải.}

Xét hàm số $y=\dfrac{3x+2}{x+1}$ trên đoạn $[0;1]$.\\

		Ta có $y'=\dfrac{1}{(x+1)^2}>0, \forall x \in [0;1]$.\\

		Suy ra hàm số đồng biến trên đoạn $[0;1]$.\\

		Do đó $m=\min\limits_{[0;1]} y=y(0)=2$; $M=\max\limits_{[0;1]} y=y(1)=\dfrac{5}{2}$.\\

		Vậy $M+m=2+\dfrac{5}{2}=\dfrac{9}{2}=4,5$.

\medskip

\end{document}
